\paragraph{}
Este capítulo detalha o protocolo experimental adotado para a avaliação das meta-heurísticas clássicas e dos sistemas multiagente propostos. Inicialmente, a \autoref{sec:soft_hard} especifica os recursos computacionais e os pacotes de software utilizados, garantindo a reprodutibilidade dos experimentos. Na sequência, a \autoref{sec:instancias} apresenta as instâncias do \textit{benchmark} de Solomon selecionadas, enquanto a \autoref{sec:implementacao} descreve os hiperparâmetros e os detalhes de implementação que regem cada estratégia de busca. Por fim, a \autoref{sec:resultados} discute os resultados obtidos sob duas perspectivas: uma análise global agregada, que utiliza métricas de \textit{gap} normalizado para comparar a eficácia das abordagens, e uma análise individualizada por instância, que ilustra a consistência da superioridade das estratégias multiagente em diversos cenários logísticos.

\section{Software e hardware}
\label{sec:soft_hard}

\paragraph{} Uma vez que o tempo de processamento é relevante para a análise deste trabalho, serão especificadas as configurações do ambiente nos quais os testes foram realizados, bem como o código fonte utilizado, permitindo a reprodutibilidade dos resultados.
\vspace{10pt}
\begin{lstlisting}[style=mystats, basicstyle=\scriptsize\ttfamily]
Sistema Operacional: Windows 11 Home 64 bits (10.0, build 22631)
Fabricante do Sistema: Dell Inc.
Modelo do Sistema: Inspiron 14 5435
BIOS: 1.17.0
Processador: AMD Ryzen 7 7730U com Radeon Graphics (16 CPUs) - 2.0GHz
Memoria: 16384MB RAM (16GB)
\end{lstlisting}

\paragraph{} O trabalho foi desenvolvido com \texttt{Python v3.11.9}, cujos principais módulos utilizados, bem como suas versões, estão listados abaixo:
\vspace{10pt}
\begin{lstlisting}[style=mystats, basicstyle=\scriptsize\ttfamily]
matplotlib              3.10.6
Mesa                    3.3.1
numpy                   2.3.3
pandas                  2.3.2
\end{lstlisting}

\paragraph{} O código fonte se encontra em: \href{https://github.com/hugodba/TCC}{https://github.com/hugodba/TCC}


\section{Instâncias de teste}
\label{sec:instancias}

\paragraph{}Neste trabalho, os experimentos numéricos são realizados utilizando um conjunto diversificado de instâncias, englobando diferentes disposições geográficas, características de janelas de tempo e capacidades de carga. Especificamente, o conjunto de testes é composto por 12 instâncias, escolhidas de forma a constituirem um conjunto representativo: \textit{c101}, \textit{c102}, \textit{c201}, \textit{c202}, \textit{r101}, \textit{r102}, \textit{r201}, \textit{r202}, \textit{rc102}, \textit{rc103}, \textit{rc201} e \textit{rc204}. Detalhes adicionais sobre os parâmetros operacionais de cada cenário podem ser verificados no \autoref{ApendiceA}.

\section{Detalhes de implementação e parametrização}
\label{sec:implementacao}

\paragraph{} Dado que a meta-heurísticas podem continuar executando por tempo indefinido, os seus hiperparâmetros foram escolhidos balizados pela performance e estagnação. Ou seja, ao longo de testes preliminares, buscou-se um resultado ótimo e o encerramento da execução o mais rápido possível após atingi-lo. Os hiperparâmetros globais de controle do MAS e do RL também buscaram resultados ótimos, porém sem a preocupação de limitar o tempo de execução quando estagnado.

\paragraph{} Para reprodutibilidade, com o intuito de mitigar a aleatoriedade inerente às meta-heurísticas, foi implementado o uso de \textit{seeds} fixas nos geradores de números aleatórios. Tal procedimento assegura que os experimentos sejam determinísticos, permitindo a comparação direta entre diferentes rodadas e configurações. Além disso, para obtenção de maior validade estatística, dada a exstência inerente de elementos estocásticos nos métodos propostos, cada algoritmo foi executado 10 vezes para cada instância e foram levados em consideração suas médias de resultado nas análises e comparações.

\paragraph{} Quanto à instrumentação dos testes, adotou-se, para mensuração do tempo, a função \texttt{time.process\_time()}, que contabiliza exclusivamente o tempo de CPU dedicado ao processo, ignorando períodos de ociosidade (\textit{sleep}) e operações de entrada/saída (I/O). Essa escolha garante que a métrica de desempenho reflita mais fielmente o esforço computacional do algoritmo.

\begin{table}[h]
\centering
\caption[\small{Parâmetros de implementação dos experimentos.}] {\label{tab:parametros} \small{Parâmetros de implementação e controle dos experimentos.}}
\footnotesize
\begin{tabular}{lll}
\toprule
\textbf{Categoria} & \textbf{Parâmetro} & \textbf{Valor} \\
\midrule
\multirow{3}{*}{\textbf{Controles}} & Número de rodadas & 10 \\
& Seed inicial & 1 \\
\midrule
\multirow{4}{*}{\textbf{MAS}} & Tamanho do \textit{Pool} & 3 \\
& Limiar de diversidade & 10 \\
& Passos por ciclo (N\_STEPS) & 20 \\
& Política de escolha & Aleatória \\
\midrule
\multirow{3}{*}{\textbf{Meta-heurísticas}} & SA (Temp. Inicial, $\alpha$, Min, Iters) & 1800, 0.99, 3, 275 \\
& TS (\textit{Max Iters}, \textit{Tenure}, \textit{Samples}) & 480, 20, 350 \\
& GA (Pop., Gen., $P_{cross}$, $P_{mut}$) & 70, 2300, 0.7, 0.05 \\
\midrule
\multirow{3}{*}{\textbf{Meta-heurísticas (MAS)}} & SA (Temp. Inicial, $\alpha$, Min, Iters) & 1800, 0.99, 3, 28 \\
& TS (\textit{Max Iters}, \textit{Tenure}, \textit{Samples}) & 80, 20, 180 \\
& GA (Pop., Gen., $P_{cross}$, $P_{mut}$) & 70, 210, 0.7, 0.05 \\
\midrule
\multirow{6}{*}{\textbf{Q-Learning}} & $\epsilon$ (Inicial / Mínimo) & 0.8 / 0.05 \\
& Decaimento de $\epsilon$ & 0.992 \\
& $\alpha$ (Taxa de aprendizado) & 0.2 \\
& $\gamma$ (Fator de desconto) & 0.8 \\
& Bin de distância & 100 \\
& Execuções por \textit{step} & 3 \\
\bottomrule
\end{tabular}
\end{table}

\section{Resultados dos experimentos}
\label{sec:resultados}

\paragraph{} Nesta seção, serão apresentados e analisados os resultados obtidos a partir dos experimentos conduzidos, com foco na comparação entre as meta-heurísticas puras e o sistema multiagente proposto. A análise será estruturada em torno de métricas de desempenho relevantes, como o número de veículos utilizados ($K$) e o tempo de execução, além de considerar a evolução das soluções ao longo do tempo e a estabilidade dos resultados. Como objetivo secundário, a distância total percorrida ($D$) também será avaliada, porém com menor ênfase. Os dados brutos de retorno dos experiments pode ser encontrado no \autoref{ApendiceC}.


\paragraph{Desempenho global}

\paragraph{}
Para possibilitar uma análise comparativa equitativa e agregada entre instâncias com ordens de grandeza de frotas e distâncias marcadamente distintas, aplicou-se um procedimento de normalização adimensional sobre os resultados brutos. A métrica adotada fundamenta-se no cálculo do \textit{gap} relativo percentual em relação às melhores soluções conhecidas na literatura (\textit{Best Known Solutions} — BKS), conforme estabelecido pelas Equações a seguir:

$$\text{Gap } K = \frac{K_{\text{médio}} - K_{\text{ref}}}{K_{\text{ref}}}$$

$$\text{Gap } D = \frac{D_{\text{médio}} - D_{\text{ref}}}{D_{\text{ref}}}$$

\paragraph{}
Onde $K_{\text{médio}}$ e $D_{\text{médio}}$ representam as médias obtidas pelo algoritmo ao longo das 10 rodadas experimentais, enquanto $K_{\text{ref}}$ e $D_{\text{ref}}$ correspondem aos valores de referência de Solomon. Sob essa ótica, um \textit{gap} igual a zero indica o alcance exato da melhor solução da literatura. Essa metodologia viabiliza o agrupamento macro de todas as 120 execuções do lote de testes em distribuições estatísticas unificadas.

\paragraph{}
O comportamento agregado dessa normalização para a frota final de veículos ($K$) e para a distância total percorrida ($D$), considerando todos os cenários combinados, pode ser observado nos diagramas de caixa apresentados na Figura~\ref{fig:boxplot_global_k} e na Figura~\ref{fig:boxplot_global_d}, respectivamente.

\begin{figure}[H]
    \centering
    \includegraphics[width=\textwidth]{K_boxplot.png}
    \caption[\small{Distribuição do gap de frota}]{\label{fig:boxplot_global_k} \small{Distribuição do gap normalizado da frota final ($K$) por método para todas as instâncias.}}
\end{figure}

\begin{figure}[H]
    \centering
    \includegraphics[width=\textwidth]{D_boxplot.png}
    \caption[\small{Distribuição do gap de distância}]{\label{fig:boxplot_global_d} \small{Distribuição do gap normalizado da distância final ($D$) por método para todas as instâncias.}}
\end{figure}

\paragraph{}
A análise macro das distribuições revela uma clara divisão de desempenho entre as abordagens tradicionais e as arquiteturas baseadas em cooperação. Ao avaliar o comportamento puramente agregado do lote completo de testes, os métodos estruturados em agência — \textit{MAS} (Sistema Multiagente) e \textit{MAS-RL} (Sistema Multiagente com Aprendizado por Reforço) — demonstram uma superioridade global evidente em relação às meta-heurísticas isoladas. Essa vantagem se traduz visualmente pelo deslocamento descendente de toda a estrutura dos diagramas de caixa, indicando uma tendência consistente de maior eficiência na minimização dos objetivos quando os solucionadores operam de forma cooperativa.

\paragraph{}
Conforme evidenciado no diagrama de frota (Figura~\ref{fig:boxplot_global_k}), a mediana do \textit{gap} de veículos para os frameworks \textit{MAS} e \textit{MAS-RL} consolida-se em patamares inferiores a 0.7, enquanto abordagens consolidadas como a Busca Tabu (\textit{TS}) e o Algoritmo Genético (\textit{GA}) operam com valores centrais próximos ou superiores a 1.0. Mais do que a redução do erro médio, as arquiteturas multiagente encolheram significativamente o intervalo interquartil da distribuição. Esse achatamento da caixa prova que o ecossistema cooperativo mitiga o risco de o algoritmo ficar aprisionado em mínimos locais medíocres ao usar as forças de cada solucionador, atingindo resultados mais consistentes, menos variáveis.

\paragraph{} No que tange à dinâmica de convergência, e não apenas ao resultado final, os métodos de agência também se destacam. A evolução temporal do \textit{gap} de frota revela que os sistemas multiagente não apenas alcançam soluções de melhor qualidade, mas o fazem de maneira mais rápida e consistente ao longo do tempo.

\paragraph{} Devido à natureza estocástica e aos tempos de execução variados entre as diferentes instâncias e rodadas, a construção dos gráficos com componente de tempo agregou os dados por meio de uma reamostragem temporal. Primeiramente, as trajetórias de cada método foram alinhadas em um \textit{grid} uniforme de 200 pontos temporais entre o instante inicial e o tempo máximo de execução, utilizando interpolação de degrau para preservar a natureza discreta das atualizações de frota.

\begin{figure}[H]
    \centering
    \includegraphics[width=\textwidth]{agregado_K.png}
    \caption[\small{Evolução temporal do gap de frota}]{\label{fig:convergencia_global_k} \small{Evolução temporal do gap de frota ($K$) para todas as instâncias combinadas.}}
\end{figure}

\paragraph{} É relevante observar que, desde o instante em que iniciam sua trajetória no gráfico, as abordagens multiagente posicionam-se em patamares de \textit{gap} significativamente inferiores aos métodos tradicionais. Esse comportamento decorre da natureza cooperativa do sistema: enquanto as meta-heurísticas puras são projetadas para uma convergência gradual ao longo de uma execução longa, os agentes do \textit{MAS} buscam resultados ótimos em cada um de seus steps, que são mais curtos. Além disso, eles partem de soluções iniciais enriquecidas promovidas pelo \textit{Path Relinking}. Desta forma, a arquitetura proposta mitiga a tradicional compensação (\textit{trade-off}) onde uma convergência rápida seria acompanhada por um desempenho final inferior, demonstrando que a integração entre agentes e memória compartilhada permite atingir resultados de alta qualidade de forma célere e sustentada.

\begin{figure}[H]
    \centering
    \includegraphics[width=\textwidth]{agregado_D.png}
    \caption[\small{Evolução temporal do gap de distância}]{\label{fig:convergencia_global_d} \small{Evolução temporal do gap de distância ($D$) para todas as instâncias combinadas.}}
\end{figure}

\paragraph{} Note que as linhas referentes a \textit{MAS} e \textit{MAS-RL} começam em um tempo maior, devido à arquitetura de implementação que recolhe as soluções parciais apenas após o término de cada \textit{step}, diferente das meta-heurísticas puras que atualizam a solução a cada iteração.


\paragraph{}
É imperativo destacar, contudo, a existência de um cenário de equivalência estatística primária ao comparar os resultados de qualidade de solução entre o \textit{MAS} puro e o \textit{MAS-RL}. Ambas as distribuições exibem caixas quase sobrepostas nos gráficos \ref{fig:boxplot_global_k} e \ref{fig:boxplot_global_d}, limites de dispersão equivalentes e medianas virtualmente idênticas no que tange ao \textit{gap} de frota e distância. A discussão da causa deste comportamento será detalhada adiante nas subseções subsequentes.

\paragraph{Desempenho por instâncias}

\paragraph{}
A análise detalhada por instância, consolidada na \autoref{fig:resultado_instancias}, reforça a superioridade das abordagens baseadas em sistemas multiagente em relação às meta-heurísticas isoladas. Observa-se que, independentemente da topologia da instância (tipo $C$, $R$ ou $RC$) ou do nível de restrição (Séries 1 e 2), as estratégias propostas demonstram um desempenho superior ou, no mínimo, equivalente às linhas de base tradicionais (\textit{SA}, \textit{GA}, \textit{TS}).

\begin{figure}[H]
    \centering
    \includegraphics[width=0.9\textwidth]{resultado_instancias.png}
    \caption[\small{Análise de convergência por instância}]{\label{fig:resultado_instancias} \small{Análise de convergência por instância. As curvas ilustram o comportamento dos métodos em diferentes topologias, evidenciando a superioridade do \textit{MAS} e \textit{MAS-RL} na exploração do espaço de busca.}}
\end{figure}


\section{Análise estatística}

\paragraph{} Um teste de hipóteses foi conduzido com o objetivo de determinar se as diferenças de desempenho observadas entre os métodos propostos possuem significância estatística. Uma vez que os dados de desempenho computacional raramente seguem uma distribuição normal em otimização estocástica, e considerando que as medições foram realizadas de forma pareada nas mesmas instâncias de Solomon, abordagens não-paramétricas foram aplicadas.

\paragraph{} Neste estudo, utilizou-se um nível de significância de $\theta = 0.05$. O procedimento estatístico foi dividido em duas etapas. Inicialmente, aplicou-se o Teste de Friedman para avaliar a existência de diferenças globais no desempenho dos cinco algoritmos testados. As hipóteses formuladas para o Teste de Friedman foram:

\begin{itemize}
    \item Hipótese Nula ($H_0$): Não existe diferença significativa entre as médias dos resultados obtidos pelos métodos avaliados; todos possuem desempenhos estatisticamente equivalentes.
    \item Hipótese Alternativa ($H_1$): Existe diferença significativa entre pelo menos dois dos algoritmos testados.
\end{itemize}

\paragraph{} O Teste de Friedman reportou uma estatística $\chi^2 = 41.48$ e um p-valor associado de $2.13 \times 10^{-8}$. Como o p-valor $\ll \theta$, a hipótese nula $H_0$ é fortemente rejeitada, confirmando que há discrepância significativa na eficácia dos algoritmos perante as instâncias analisadas. Os resultados agregados que basearam este teste estão sumarizados na Tabela \ref{tab:resultados_agregados}.

\begin{table}[htbp]
    \centering
    \caption[\small{Matriz de Resultados Agregados}]{\label{tab:resultados_agregados} \small{Matriz de Resultados Agregados (Média de Veículos por Instância)}}
    \begin{tabularx}{\textwidth}{l *{5}{>{\centering\arraybackslash}X}}
        \toprule
        \textbf{Instância} & \textbf{SA} & \textbf{TS} & \textbf{GA} & \textbf{MAS} & \textbf{MAS-RL} \\
        \midrule
        \texttt{c101}  & 18.2 & 18.8 & 19.1 & 17.0 & 17.1 \\
        \texttt{c102}  & 15.2 & 17.0 & 16.5 & 15.0 & 14.8 \\
        \texttt{c201}  & 11.1 & 12.1 & 11.3 & 10.1 & 10.2 \\
        \texttt{c202}  & 8.5  & 10.6 & 9.2  & 7.6  & 7.7  \\
        \texttt{r101}  & 25.2 & 25.5 & 26.7 & 24.0 & 24.0 \\
        \texttt{r102}  & 21.1 & 21.6 & 22.5 & 20.1 & 20.5 \\
        \texttt{r201}  & 9.3  & 12.5 & 9.8  & 9.1  & 9.0  \\
        \texttt{r202}  & 7.8  & 11.1 & 7.8  & 7.2  & 7.4  \\
        \texttt{rc102} & 19.0 & 20.4 & 21.4 & 18.1 & 18.0 \\
        \texttt{rc103} & 16.3 & 17.7 & 18.7 & 16.3 & 16.2 \\
        \texttt{rc201} & 10.5 & 13.2 & 10.5 & 9.5  & 9.4  \\
        \texttt{rc204} & 5.2  & 5.8  & 5.0  & 4.6  & 4.6  \\
        \bottomrule
    \end{tabularx}
\end{table}

\paragraph{} Comprovada a diferença global, procedeu-se à segunda etapa da análise através do Teste de Nemenyi, um teste \textit{post-hoc} utilizado para identificar empiricamente, par a par, quais algoritmos apresentam diferenças significativas entre si. A Tabela \ref{tab:nemenyi_pvalues} sumariza os \textit{ranks} médios atribuídos a cada algoritmo, onde menores valores indicam maior aderência aos melhores resultados de otimização, bem como a matriz de p-valores das comparações pareadas.

\begin{table}[htbp]
    \centering
    \caption[\small{Matriz de p-valores do Teste Post-Hoc de Nemenyi e ranks médios}]{\label{tab:nemenyi_pvalues} \small{Matriz de p-valores do Teste Post-Hoc de Nemenyi e ranks médios}}
    \begin{tabular}{lcccccc}
        \toprule
        \multirow{2}{*}{\textbf{Método}} & \multirow{2}{*}{\textbf{Rank Médio}} & \multicolumn{5}{c}{\textbf{p-valor vs. Algoritmo}} \\
        \cmidrule(lr){3-7}
        & & \textbf{MAS-RL} & \textbf{MAS} & \textbf{SA} & \textbf{GA} & \textbf{TS} \\
        \midrule
        \textbf{MAS-RL} & \textbf{1.50} & -      & 1.0000 & 0.0867 & 0.0002 & 0.0000 \\
        \textbf{MAS}     & \textbf{1.54} & 1.0000 & -      & 0.1016 & 0.0003 & 0.0000 \\
        \textbf{SA}      & 3.12 & 0.0867 & 0.1016 & -      & 0.4076 & 0.1583 \\
        \textbf{GA}      & 4.25 & 0.0002 & 0.0003 & 0.4076 & -      & 0.9858 \\
        \textbf{TS}      & 4.58 & 0.0000 & 0.0000 & 0.1583 & 0.9858 & -      \\
        \bottomrule
    \end{tabular}
\end{table}

\paragraph{} Os testes indicam uma vantagem estatística notória dos sistemas multiagentes. O modelo \textit{MAS-RL} atingiu o menor \textit{rank} médio ($1.50$), intimamente seguido pelo modelo cooperativo puro MAS ($1.54$). Em contraste, as meta-heurísticas isoladas obtiveram ranks substancialmente maiores (SA: $3.12$; GA: $4.25$; TS: $4.58$).

\paragraph{} A matriz de p-valores gerada pelo Teste de Nemenyi (Tabela \ref{tab:nemenyi_pvalues}) comprova que o \textit{MAS-RL} supera significativamente tanto o TS (p-valor = $0.0000$) quanto o GA (p-valor = $0.0002$). Destaca-se também que, para a métrica de veículos, a diferença entre o \textit{MAS-RL} e o SA (p-valor = $0.0867$) encontra-se no limite tendencial, não rejeitando $H_0$ estritamente sob $\theta = 0.05$. Por fim, os resultados do \textit{MAS} e do \textit{MAS-RL} (p-valor = $1.0000$) são considerados estatisticamente equivalentes, ressaltando o ganho de robustez proveniente da arquitetura do Sistema Multiagente em si.

\paragraph{} Para ilustrar estas conclusões de forma visual, a Figura \ref{fig:cd_diagram} apresenta o Diagrama de Diferença Crítica (\textit{Critical Difference Diagram}). Os algoritmos não conectados por barras horizontais possuem diferenças estatísticas significativas de desempenho, atestando a superioridade dos arranjos cooperativos.

\begin{figure}[htbp]
    \centering
    \includegraphics[width=0.85\textwidth]{cd_diagram.png}
    \caption[\small{Diagrama de Diferença Crítica (Nemenyi) para a minimização de Veículos.}]{\label{fig:cd_diagram} \small{Diagrama de Diferença Crítica (Nemenyi) para a minimização de Veículos.}}
\end{figure}

\paragraph{} Esses resultados acarretam importantes implicações práticas para o modelo desenvolvido. A superioridade estatisticamente provada do \textit{MAS} e do \textit{MAS-RL} confirma que a arquitetura cooperativa mitiga o risco de paralisação em ótimos locais, gerando configurações de frota consideravelmente melhores nas instâncias de Solomon sem que esse ganho seja decorrente do acaso. Em relação ao \textit{SA}, embora a diferença de desempenho seja evidente, a ausência de significância estatística estrita sugere que, para algumas instâncias, o SA pode alcançar resultados próximos aos sistemas multiagente, reforçando a necessidade de análises mais granulares por cenário.